\documentclass[11pt]{amsart}
\usepackage{amsmath,amssymb,amsthm}
\usepackage[margin=1.05in]{geometry}

\numberwithin{equation}{section}
\newtheorem{theorem}{Theorem}[section]
\newtheorem{proposition}[theorem]{Proposition}
\newtheorem{lemma}[theorem]{Lemma}
\newtheorem{corollary}[theorem]{Corollary}
\theoremstyle{definition}
\newtheorem{definition}[theorem]{Definition}
\theoremstyle{remark}
\newtheorem{remark}[theorem]{Remark}
\newtheorem{example}[theorem]{Example}

\newcommand{\Z}{\mathbb{Z}}
\newcommand{\Q}{\mathbb{Q}}
\newcommand{\R}{\mathbb{R}}
\newcommand{\N}{\mathbb{N}}
\newcommand{\C}{\mathbb{C}}
\newcommand{\T}{\mathbb{T}}
\newcommand{\cA}{\mathcal{A}}
\newcommand{\cB}{\mathcal{B}}
\newcommand{\cG}{\mathcal{G}}
\newcommand{\cL}{\mathcal{L}}
\newcommand{\cR}{\mathcal{R}}
\newcommand{\cU}{\mathcal{U}}
\newcommand{\pp}{\mathfrak{p}}
\newcommand{\aaa}{\mathfrak{a}}
\newcommand{\qq}{\mathfrak{q}}
\newcommand{\hO}{\widehat{\mathcal{O}}}
\newcommand{\Ns}{\mathrm{N}}
\newcommand{\Cl}{\operatorname{Cl}}
\newcommand{\Ent}{\mathrm{S}}
\newcommand{\Hom}{\operatorname{Hom}}
\newcommand{\coker}{\operatorname{coker}}
\newcommand{\Aut}{\operatorname{Aut}}
\newcommand{\htop}{\mathrm{ht}}
\newcommand{\id}{\mathrm{id}}
\newcommand{\Msym}{M_{\mathrm{sym}}}
\newcommand{\phisym}{\varphi_{\mathrm{sym}}}
\newcommand{\nusym}{\nu_{\mathrm{sym}}}
\newcommand{\Prob}{\operatorname{Prob}}

\PassOptionsToPackage{hyphens}{url}
\usepackage[hidelinks,bookmarks=false]{hyperref}

\title[Invariants, entropy, and the measure class]{Localized Bost--Connes systems III:\\ invariants, entropy, and the measure class}
\author{Ruqing Chen}
\address{GUT Geoservice Inc., Montreal, Canada}
\email{ruqing@hotmail.com}
\thanks{Sources, code and data for the entire series: \url{https://github.com/Ruqing1963/localized-bost-connes}; this paper is in \texttt{papers/III-invariants/}.}
\date{}

\begin{document}

\begin{abstract}
The obstruction group $\Xi_\beta(K,S)$ of \cite{Main} is a complete invariant of the phase
structure of the localized Bost--Connes system. We determine what kind of invariant it is,
and the answer organizes everything else: \emph{$\Xi_\beta$ is an invariant of a measure
class, not of a $C^*$-algebra}. It is the kernel of the natural map from $\widehat{G_S}$ to
the \emph{measurable} first cohomology of the tail equivalence relation taken with respect
to $\pi_\beta$, and the whole of its dependence on $\beta$ sits in that measure class.

Three consequences follow, and they explain a series of otherwise puzzling vanishings.
First, at the level of $K$-theory the KMS simplex is exactly the set of states $\omega$ on
the pointed ordered group $K_0(C(Y_S))=C(Y_S,\Z)$ satisfying the eigenvalue equations
$\omega\circ\alpha_{\pp*}=\Ns\pp^{-\beta}\omega$; and since the image of
$K_0(C(Y_S))$ in $K_0(\cA_{K,S})$ is the group of $\alpha_*$-coinvariants, the passage to
the crossed product annihilates precisely those eigenvectors. Second, every homotopy-invariant construction is blind:
Connes--Thom gives $K_0(\cB_{K,S})=0$ for the continuous core, so every dual-trace pairing
on it vanishes identically; the dual flow acts trivially on $K_*(\cB_{K,S})$ because $\R$ is
connected; and the groupoid homology of $\cG_S$, with or without twists by the modular
cocycle, carries no $\beta$ at all. Third, no dynamical entropy is available:
$\alpha_\aaa$ and the transfer operator $\cL_\aaa$ preserve no $\mathrm{KMS}_\beta$ state
and indeed no invariant probability measure exists, in accordance with the type
$\mathrm{III}$ classification; Voiculescu's topological entropy of $\alpha_\aaa$ vanishes
because multiplication by $\aaa$ is an ultrametric contraction; and the modular flow, which
does preserve every KMS state, has zero CNT entropy because it lies in the compact gauge
group.

What does see the phase structure is relative. We prove
$\Ent(\omega_{\varphi_\varepsilon}\|\omega_{\phisym})=|A_\beta|\log2$, exactly $\log2$ per
threshold crossed and $+\infty$ once infinitely many have been, and show that the entire
value is carried by the centre, so the type $\mathrm{III}$ factors contribute nothing. The
resulting staircase coincides with the subfactor index staircase of Paper II of the series.
We close by correcting the thermodynamic reading: all $\mathrm{KMS}_\beta$ states at a given
$\beta$ are equilibrium states on an equal footing, so the condensation of a phase bit is
not a Landauer erasure.

Along the way we fix two structural boundaries that are easy to misplace. A
$\mathrm{KMS}_\beta$ state with $\beta\ne0$ is not a trace, so it defines no functional on
$K_0(\cA_{K,S})$; the $K$-theoretic locus of the KMS states is the eigenvalue problem on
$K_0(C(Y_S))$ instead. And the de la Harpe--Skandalis determinant is natively defined on
$\cU_0$, which is the subgroup by which $K_1=\cU/\cU_0$ is the quotient, so it induces no
invariant on $K_1$.
\end{abstract}

\maketitle

\tableofcontents

\section{Introduction: which kind of invariant}

We use the notation of \cite{Main}: $\cA_{K,S}=C(Y_S)\rtimes J_S$ with
$\sigma_t(\mu_\aaa)=\Ns\aaa^{it}\mu_\aaa$, symmetry group $G_S$, and
\begin{equation}\label{eq:Xi}
\Xi_\beta(K,S)=\Bigl\{\chi\in\widehat{G_S}:\ \sum_{\pp\in S}\Ns\pp^{-\beta}\bigl|1-\chi(\sigma_\pp)\bigr|^2<\infty\Bigr\},
\end{equation}
the $\mathrm{KMS}_\beta$ simplex being $\Prob(G_S/\Xi_\beta^\perp)$. Write
$\cR_S$ for the tail equivalence relation on $V=\Z_{\ge0}^S$ and $\pi_\beta$ for the forced
product measure of \cite[Lem.~2.10]{Main}, $\alpha_\pp$ for the endomorphism
$f\mapsto\mu_\pp f\mu_\pp^*$ of $C(Y_S)$, and $\cB_{K,S}=\cA_{K,S}\rtimes_\sigma\R$ for the
continuous core.

\emph{The series.} The present paper is III in the series. The foundational paper is
\cite{Main}, which classifies the $\mathrm{KMS}$ states and introduces $\Xi_\beta$; Paper I
\cite{SeriesI} treats factor type and the spectrum of transitions. Papers II (ideal
structure, boundary quotient, index), IV (noncommutative metric geometry) and V (the free
variant) are in preparation, and nothing here depends on them; we refer to them by numeral,
without citation.

The organizing statement is Theorem~\ref{thm:meas}: $\Xi_\beta$ is the kernel of a map into
a \emph{measurable} cohomology group computed relative to $\pi_\beta$. Since a
$C^*$-algebra, a groupoid and a homotopy functor all forget a measure class, one should
expect them to forget $\Xi_\beta$, and \S\S\ref{sec:Ktheory}--\ref{sec:entropy} confirm this
in each case. The point of collecting the vanishings is not that they are surprising
individually --- several are easy --- but that they have a single cause, and that the cause
identifies where the invariant does live.

\section{$\Xi_\beta$ is measurable first cohomology}\label{sec:meas}

Let $\sigma:\Z^{(S)}\to G_S$ be the canonical homomorphism, regarded as a cocycle on
$\cR_S$ via $\sigma(v+a,v)=\sigma_\aaa$.

\begin{lemma}[Coboundary dictionary]\label{lem:dict}
Let $A$ be a second countable locally compact abelian group, $\Phi:\Z^{(S)}\to A$ a
homomorphism, $\psi\in\widehat A$, and $c_\pp=\overline{\psi(\Phi(e_\pp))}\in\T$. The
following are equivalent.
\begin{enumerate}
\item $\psi\circ\Phi$ is a coboundary on $(V,\pi_\beta,\cR_S)$: there is measurable
$g:V\to\T$ with $g(v+e_\pp)=c_\pp g(v)$ a.e., for every $\pp$;
\item $\sum_\pp\bigl(1-|m_\pp|\bigr)<\infty$, where $m_\pp=(1-t_\pp)/(1-c_\pp t_\pp)$ and
$t_\pp=\Ns\pp^{-\beta}$;
\item $\sum_\pp\Ns\pp^{-\beta}\bigl|1-\psi(\Phi(e_\pp))\bigr|^2<\infty$.
\end{enumerate}
Moreover, if these hold, then $\psi$ is trivial on the group $E(\Phi)$ of essential values
of $\Phi$; so $E(\Phi)\subseteq\Lambda^\perp$ with
$\Lambda=\{\psi:\psi\circ\Phi\ \text{a coboundary}\}$.
\end{lemma}

\begin{proof}
$(1)\Leftrightarrow(2)$ is Propositions 3.3 and 3.4 of \cite{Main}; those proofs use only
that $c_\pp\in\T$, that the coordinates $v_\pp$ are independent under $\pi_\beta$ with
$\mu_\pp(k)=(1-t_\pp)t_\pp^k$, and that $|g|=1$. The target group of $\psi$ plays no role.
$(2)\Leftrightarrow(3)$ is \cite[Lem.~3.2]{Main} together with
$|1-c_\pp|^2=2(1-\operatorname{Re}c_\pp)$ and $|1-c_\pp|=|1-\overline{c_\pp}|$.
For the last statement, a coboundary has no essential value other than $1$, and
$\psi$ maps essential values of $\Phi$ to essential values of $\psi\circ\Phi$
\cite{Schmidt}; so $\psi(E(\Phi))=\{1\}$.
\end{proof}

\begin{remark}\label{rem:noconverse}
The converse, that $\psi$ trivial on $E(\Phi)$ makes $\psi\circ\Phi$ a coboundary, holds
when $\Phi$ is regular in the sense of Schmidt \cite{Schmidt}, i.e.\ cohomologous to a
cocycle with values in $E(\Phi)$, and then $E(\Phi)=\Lambda^\perp$ by Pontryagin duality.
For $A=G_S$ compact this is automatic, and Theorem~\ref{thm:meas} needs only
$(1)\Leftrightarrow(3)$. For the cocycle with values in $G_S\times\R$ that governs the
factor type, regularity is not known, and \cite[Prop.~2.2 and Rem.~2.3]{SeriesI} obtains
from the inclusion $E(\Phi)\subseteq\Lambda^\perp$ an upper bound on the type only.
\end{remark}

\begin{theorem}[Cohomological description]\label{thm:meas}
For every $\beta>0$,
\begin{equation}\label{eq:coh}
\Xi_\beta(K,S)\ =\ \ker\Bigl(\widehat{G_S}\ \xrightarrow{\ \chi\mapsto[\chi\circ\sigma]\ }\ H^1_{\mathrm{meas}}\bigl(\cR_S,\pi_\beta;\T\bigr)\Bigr),
\end{equation}
measurable $1$-cohomology relative to the measure class of $\pi_\beta$, in the sense of
Feldman--Moore \cite{FM}.
\end{theorem}

\begin{proof}
Apply Lemma~\ref{lem:dict} with $A=G_S$, $\Phi=\sigma$: condition (3) is \eqref{eq:Xi}, and
(1) says the class of $\chi\circ\sigma$ vanishes in $H^1_{\mathrm{meas}}$.
\end{proof}

\begin{corollary}[Where the $\beta$-dependence sits]\label{cor:where}
In \eqref{eq:coh} the objects $\widehat{G_S}$, $\cR_S$ and $\sigma$ are $\beta$-independent.
The entire dependence on $\beta$ is carried by the measure class of $\pi_\beta$, through
which $H^1_{\mathrm{meas}}$ is computed.
\end{corollary}

\begin{remark}\label{rem:category}
Corollary~\ref{cor:where} is the reason for everything below. $K$-theory is an invariant of
a $C^*$-algebra; groupoid homology is an invariant of a topological groupoid; dynamical
entropy is an invariant of a measure-preserving map. None of the three retains a measure
class that is merely quasi-invariant. The vanishings of
\S\S\ref{sec:Ktheory}--\ref{sec:entropy} are therefore not failures of ingenuity but
category errors, and we record them in that spirit.
\end{remark}

\section{$K$-theory: the eigenvalue equation}\label{sec:Ktheory}

We first correct a statement announced in an earlier draft.

\begin{proposition}[KMS states do not act on $K_0(\cA_{K,S})$]\label{prop:notrace}
Let $\varphi$ be a $\mathrm{KMS}_\beta$ state with $\beta>0$ and $S\ne\emptyset$. Then
$[p]\mapsto\varphi(p)$ does \emph{not} define a homomorphism $K_0(\cA_{K,S})\to\R$.
\end{proposition}

\begin{proof}
A positive functional descends to $K_0$ only if it is constant on Murray--von Neumann
classes, i.e.\ only if $\varphi(vv^*)=\varphi(v^*v)$ for all $v$, i.e.\ only if $\varphi$ is
a trace. Take $v=\mu_\pp$: then $v^*v=1$ and $vv^*=\mathbf 1_{\pp Y_S}$, so
$[1]=[\mathbf 1_{\pp Y_S}]$ in $K_0$, while the KMS condition gives $\varphi(1)=1$ and
$\varphi(\mathbf 1_{\pp Y_S})=\Ns\pp^{-\beta}<1$.
\end{proof}

The remedy is to work on the coefficient algebra, where a KMS state restricts to a measure
and hence to a trace.

\begin{theorem}[The KMS simplex as an eigenvalue problem]\label{thm:eigen}
Let $\alpha_{\pp*}$ be the endomorphism $\mathbf 1_U\mapsto\mathbf 1_{\pp U}$ of the pointed
ordered group $\bigl(K_0(C(Y_S)),K_0^+,[1]\bigr)=\bigl(C(Y_S,\Z),C(Y_S,\Z)_+,1\bigr)$. Then
$\varphi\mapsto\omega_\varphi$, $\omega_\varphi[\mathbf 1_U]:=\varphi(\mathbf 1_U)$, is an
affine bijection
\[
\mathrm{KMS}_\beta\bigl(\cA_{K,S}\bigr)\ \xrightarrow{\ \sim\ }\
\Bigl\{\omega\ \text{a state of }\bigl(C(Y_S,\Z),C(Y_S,\Z)_+,1\bigr):\
\omega\circ\alpha_{\pp*}=\Ns\pp^{-\beta}\,\omega\ \ \forall\pp\in S\Bigr\}.
\]
\end{theorem}

\begin{proof}
By \cite[Prop.~2.8]{Main} a $\mathrm{KMS}_\beta$ state is $\nu\circ E$ for a unique
scaling-invariant $\nu$, and $\nu$ is a Borel probability measure on the compact totally
disconnected space $Y_S$, hence determined by its values on clopen sets; conversely a
positive normalized finitely additive assignment on the clopen algebra extends to a regular
Borel measure. Since $C(Y_S)$ is commutative, $\nu$ is a trace and $\nu_*$ is a well-defined
state of the pointed ordered group. The scaling relation $\nu(\aaa E)=\Ns\aaa^{-\beta}\nu(E)$
reads on clopen sets as $\omega(\alpha_{\pp*}[\mathbf 1_U])=\Ns\pp^{-\beta}\omega[\mathbf 1_U]$,
and the relations for $\pp\in S$ generate those for all $\aaa\in J_S$.
\end{proof}

\begin{corollary}[The crossed product annihilates the datum]\label{cor:annih}
In $K_0(\cA_{K,S})$ one has $\iota_*[\mathbf 1_U]=\iota_*[\mathbf 1_{\pp U}]$ for every
clopen $U\subseteq Y_S$ and $\pp\in S$, since $\mu_\pp\mathbf 1_U$ is a partial isometry from
$\mathbf 1_U$ to $\mathbf 1_{\pp U}$. So $\iota_*:C(Y_S,\Z)\to K_0(\cA_{K,S})$ factors
through the group of coinvariants $C(Y_S,\Z)/\sum_\pp(1-\alpha_{\pp*})C(Y_S,\Z)$, and a
functional on $C(Y_S,\Z)$ that factors through $K_0(\cA_{K,S})$ satisfies
$\omega\circ\alpha_{\pp*}=\omega$ for every $\pp$. By Theorem~\ref{thm:eigen} the
$\mathrm{KMS}_\beta$ functionals satisfy instead $\omega\circ\alpha_{\pp*}=\Ns\pp^{-\beta}\omega$
with $\Ns\pp^{-\beta}\ne1$ for $\beta\ne0$. Hence no nonzero $\mathrm{KMS}_\beta$ functional
factors through $K_0(\cA_{K,S})$: passing from the coefficient algebra to the crossed
product is exactly the quotient that kills the eigenvectors of interest. (For a single
prime the image of $\iota_*$ is the whole cokernel $\coker(1-\alpha_*)=K_0(\cA_{K,S})$ and
$K_1=0$; for several primes the image is the bottom piece of a filtration of $K_0$ by Koszul
homology, computed in Paper II of the series.)
\end{corollary}

\begin{remark}\label{rem:pair}
Theorem~\ref{thm:eigen} is a statement about the \emph{pair} $\bigl(K_0(C(Y_S)),\alpha_*\bigr)$
together with a real parameter, not about a single abelian group. Every functor that forgets
$\alpha_*$, or that forces it to act as the identity, loses the arithmetic; and
$K_0$ of the crossed product does the latter by construction.
\end{remark}

\section{Homotopy invariants are blind}\label{sec:homotopy}

\begin{theorem}[The continuous core]\label{thm:core}
For every $K$, $S$ and $\beta>0$:
\begin{enumerate}
\item $K_0(\cB_{K,S})\cong K_1(\cA_{K,S})$ and $K_1(\cB_{K,S})\cong K_0(\cA_{K,S})$;
\item the dual flow $\hat\sigma$ induces the identity on $K_*(\cB_{K,S})$;
\item the dual weight $\hat\varphi$ of any $\mathrm{KMS}_\beta$ state, a faithful
semifinite trace on $\cB_{K,S}$, vanishes on every projection in its domain (in every
matrix algebra over $\cB_{K,S}$); in particular the induced map $\hat\varphi_*$ on $K_0$ of
its domain ideal is zero, and the pairing of $\hat\varphi$ with $K_0(\cB_{K,S})$ carries no
information.
\end{enumerate}
Consequently no invariant built from $\bigl(K_*(\cB_{K,S}),\hat\sigma_*,\text{trace pairings}\bigr)$
distinguishes systems with different $\Xi_\beta$.
\end{theorem}

\begin{proof}
(1) is the Connes--Thom isomorphism $K_i(\cA\rtimes_\sigma\R)\cong K_{i+1}(\cA)$ \cite{CT}.

(2) The map $s\mapsto\hat\sigma_s$ is a point-norm continuous path of $*$-automorphisms with
$\hat\sigma_0=\id$, hence a homotopy of $*$-homomorphisms $\cB\to\cB$; $K$-theory is homotopy
invariant, so $\hat\sigma_{s*}=\id$ for every $s$.

(3) The defining property of the dual weight of a $\mathrm{KMS}_\beta$ state is that it is a
trace scaled by the dual action, $\hat\varphi\circ\hat\sigma_s=e^{-\beta s}\hat\varphi$;
in particular its domain is $\hat\sigma$-invariant. Let $p$ be a projection in the domain.
Then $s\mapsto\hat\sigma_s(p)$ is a norm-continuous path of projections in the domain, so
$\hat\sigma_s(p)=u_spu_s^*$ for a unitary $u_s$ of the unitization, and the trace property
$\hat\varphi(x^*x)=\hat\varphi(xx^*)$ gives $\hat\varphi(\hat\sigma_s(p))=\hat\varphi(p)$. Hence
$\hat\varphi(p)=e^{-\beta s}\hat\varphi(p)$ for every $s$, and $\hat\varphi(p)=0$.
\end{proof}

\begin{remark}[Two mechanisms, one cause]\label{rem:twomech}
Part (3) is proved here without computing $K_0(\cB_{K,S})$, and this is deliberate. By (1),
$K_0(\cB_{K,S})\cong K_1(\cA_{K,S})$. For $|S|=1$ this is $\ker(1-\alpha_*)$ on $C(Y_S,\Z)$
by Pimsner's sequence, and it vanishes: a fixed point of $\alpha_*$ is supported in
$\pp^nY_S$ for every $n$, hence in the nowhere dense set $\partial Y_S=\bigcap_n\pp^nY_S$,
hence is $0$. For $|S|\ge2$ the group $K_1(\cA_{K,S})$ is built from the odd Koszul homology
of the commuting family $(1-\alpha_{\pp*})$, and we have not determined it (Paper II). Were
it zero, (3) would follow trivially. The argument above shows that (3) holds regardless: the
scaling $\hat\varphi\circ\hat\sigma_s=e^{-s}\hat\varphi$ and the homotopy triviality of
$\hat\sigma_*$ are incompatible with a nonzero pairing, whatever the group is. It is the
\emph{connectedness of $\R$} that kills the pairing, not the vanishing of a $K$-group ---
the same mechanism as in Remark~\ref{rem:connected}.
\end{remark}

\begin{remark}[Connected versus totally disconnected]\label{rem:connected}
The obstruction in (2) is general: any action of a \emph{connected} group induces the
trivial action on $K$-theory. The modular flow is an action of $\R$, so the scaling
structure it carries (namely $\hat\varphi\circ\hat\sigma_s=e^{-s}\hat\varphi$, where the type
and the flow of weights live) is invisible to $K$-theory in principle. By contrast $G_S$ is
profinite, hence totally disconnected, and acts nontrivially on $K_0(C(Y_S))=C(Y_S,\Z)$ by
permuting clopen sets. Arithmetic symmetry, being Galois-theoretic and so profinite, can
register on $K$-theory; modular dynamics, being a flow, cannot.
\end{remark}

\begin{proposition}[The determinant does not act on $K_1$]\label{prop:dlhs}
Let $A$ be unital with a trace $\tau$. The de la Harpe--Skandalis determinant \cite{dlHS} is
a homomorphism $\Delta_\tau:\cU_0(A)\to\R/\tau_*(K_0(A))$ on the connected component of the
unitary group. Since $K_1(A)=\cU_\infty(A)/\cU_{\infty,0}(A)$, the domain of $\Delta_\tau$ is
precisely the subgroup by which $K_1$ is the quotient; $\Delta_\tau$ descends to
$\cU_0/\overline{D\cU_0}$ and does not induce a map $K_1(A)\to\R$.
\end{proposition}

\begin{proof}
Definitional: $\Delta_\tau$ vanishes on commutators of $\cU_0$, and a homomorphism defined
on $\cU_0$ carries no information about $\cU/\cU_0$.
\end{proof}

\begin{remark}[Parity]\label{rem:parity}
A trace is a cyclic $0$-cocycle and pairs with $K_0$; pairing with $K_1$ requires a cyclic
$1$-cocycle. On the core the natural candidate is
$\psi(b_0,b_1)=\hat\varphi(b_0\hat\delta(b_1))$ with $\hat\delta$ the generator of the dual
flow, but cyclicity needs $\hat\varphi\circ\hat\delta=0$ whereas the defining property of the
dual weight is $\hat\varphi\circ\hat\sigma_s=e^{-s}\hat\varphi$, i.e.\
$\hat\varphi\circ\hat\delta=-\hat\varphi\ne0$. The pair $(\hat\varphi,\hat\delta)$ produces a
\emph{twisted} cyclic $1$-cocycle, consistently with Theorem~\ref{thm:meas}.
\end{remark}

\begin{proposition}[Groupoid homology carries no $\beta$]\label{prop:homology}
$H_*(\cG_S;\Z)\cong H_*\bigl(I_S;C_c(\widetilde Y_S,\Z)\bigr)$ is defined without reference
to $\beta$, as is any twist by the modular cocycle $c(\aaa\mathfrak b^{-1})=\log\Ns\aaa-\log\Ns\mathfrak b$,
which does not involve $\beta$. Hence neither can determine $\Xi_\beta$, which varies with
$\beta$. Twisting the coefficients instead does not help: with $\R_\beta$ denoting $\R$ with
$I_S$ acting by $\aaa\mapsto\Ns\aaa^{-\beta}$,
\[
H^n\bigl(I_S;\R_\beta\bigr)=0\quad\text{for all }n\ge0,
\]
while $H^0$ with coefficients in the measures twisted by $e^{-\beta c}$ returns the KMS
simplex tautologically.
\end{proposition}

\begin{proof}
The first statements are immediate: $\beta$ enters only as the exponent in the
Radon--Nikodym derivative, i.e.\ in the measure class. For the vanishing, pick $\pp\in S$;
the generator $\pp$ acts on $\R_\beta$ by $\lambda=\Ns\pp^{-\beta}\ne1$, and for $\Z$ acting
on $\R$ by $\lambda\ne1$ one has $H^0=\ker(\lambda-1)=0$ and $H^1=\R/(\lambda-1)\R=0$; the
Koszul complex for $I_S=\Z^n$ is then exact, since one of its generators acts by an
automorphism $\lambda-1$ of the coefficients, and for infinite $S$ the Milnor sequence of
the increasing union of finite-rank subgroups gives the same vanishing. The last statement is the definition of
quasi-invariance together with \cite[Prop.~2.8]{Main}.
\end{proof}

\section{Dynamical entropy}\label{sec:entropy}

\begin{proposition}[Scaling, not invariance]\label{prop:scaling}
For a $\mathrm{KMS}_\beta$ state $\varphi$, $\aaa\in J_S$ and $x\in\cA_{K,S}$,
\[
\varphi\bigl(\alpha_\aaa(x)\bigr)=\Ns\aaa^{-\beta}\varphi(x),
\qquad
\varphi\bigl(\cL_\aaa(x)\bigr)=\Ns\aaa^{\beta}\,\varphi\bigl(x\,\mathbf 1_{\aaa Y_S}\bigr),
\]
where $\cL_\aaa(x)=\mu_\aaa^*x\mu_\aaa$ is the unital completely positive transfer operator.
Neither map preserves $\varphi$ for $\aaa\ne(1)$, $\beta>0$.
\end{proposition}

\begin{proof}
$\sigma_{i\beta}(\mu_\aaa)=\Ns\aaa^{-\beta}\mu_\aaa$ and
$\sigma_{i\beta}(\mu_\aaa^*)=\Ns\aaa^{\beta}\mu_\aaa^*$. The KMS condition
$\varphi(uv)=\varphi(v\sigma_{i\beta}(u))$ with $u=\mu_\aaa$, $v=x\mu_\aaa^*$ gives the
first, using $\mu_\aaa^*\mu_\aaa=1$; with $u=\mu_\aaa^*$, $v=x\mu_\aaa$ it gives the second.
For the failure of the second, take $x=\mathbf 1_E$ with $E=Y_S\setminus\aaa Y_S$: the right
side is $0$ while $\varphi(\mathbf 1_E)=1-\Ns\aaa^{-\beta}>0$.
\end{proof}

\begin{theorem}[No invariant measure, hence no metric entropy]\label{thm:noinv}
For $S\ne\emptyset$ there is no $I_S$-invariant Borel probability measure on
$\widetilde Y_S$ other than those concentrated on the finite boundary $\partial Y_S$, and no
such measure is a $\mathrm{KMS}_\beta$ state for $\beta>0$. Consequently the
Connes--St\o rmer, CNT and Sauvageot--Thouvenot entropies \cite{CS,CNT,ST}
$h_\varphi(\alpha_\aaa)$ and $h_\varphi(\cL_\aaa)$ are undefined for every
$\mathrm{KMS}_\beta$ state.
\end{theorem}

\begin{proof}
An invariant $\nu$ satisfies $\nu(\aaa^{-1}Y_S)=\nu(Y_S)$ for all $\aaa\in J_S$, and
$\widetilde Y_S=\bigcup_\aaa\aaa^{-1}Y_S$, so $\nu(Y_S)=1$; then $\nu(\pp^kY_S)=1$ for all $k$
and $\nu$ is concentrated on $\bigcap_k\pp^kY_S$; intersecting over $\pp\in S$ concentrates
$\nu$ on $\partial Y_S=\{[0,g]\}\cong G_S/r(U_S)\cong\Cl^+_S$ \cite[(2.1)]{Main}, a finite
set. A
$\mathrm{KMS}_\beta$ state has $\varphi(\mathbf 1_{\pp Y_S})=\Ns\pp^{-\beta}<1$ and so is not
of this form. All three entropy definitions require the state to be invariant under the map
whose entropy is computed; Proposition~\ref{prop:scaling} shows no KMS state is.
\end{proof}

\begin{remark}[Consistency with type $\mathrm{III}$]\label{rem:typeIII}
Theorem~\ref{thm:noinv} restates the type classification. A type $\mathrm{III}$ factor has no
semifinite trace, the associated relation has no invariant measure in the measure class, and
Kolmogorov--Sinai theory is unavailable by construction.
\end{remark}

\begin{theorem}[Topological entropy vanishes]\label{thm:ht}
Let $T_\aaa:Y_S\to Y_S$, $[x,g]\mapsto[\aaa x,\sigma_\aaa g]$, so that
$\cL_\aaa(f)=\mu_\aaa^*f\mu_\aaa=f\circ T_\aaa$ on $C(Y_S)$. Then $h_{\mathrm{top}}(T_\aaa)=0$.
On the quotient of $\cA_{K,S}$ in which every $\mu_\pp$ is unitary (the boundary quotient
of Paper II) the map $\alpha_\aaa$ becomes the inner automorphism $\mathrm{Ad}\,\mu_\aaa$,
whose Voiculescu--Brown topological entropy \cite{Voiculescu,Brown} is $0$.
\end{theorem}

\begin{proof}
$Y_S$ is the quotient of $\hO_S\times G_S$ by $U_S$, and $T_\aaa$ lifts to
$\widetilde T_\aaa(x,g)=(\aaa x,\sigma_\aaa g)$. Put
\[
d\bigl((x,g),(x',g')\bigr)=\sup_\pp c_\pp\Ns\pp^{-v_\pp(x_\pp-x'_\pp)}+d_G(g,g'),
\]
with $c_\pp$ summable and $d_G$ an invariant metric on $G_S$; one has $v_\pp(\aaa(x-x'))=v_\pp(\aaa)+v_\pp(x-x')$ and
$d_G(\sigma_\aaa g,\sigma_\aaa g')=d_G(g,g')$, so $\widetilde T_\aaa$ is $1$-Lipschitz.
Hence the Bowen metrics satisfy $d_n=d$, the $(n,\epsilon)$-separated sets are the
$\epsilon$-separated sets, their cardinality is independent of $n$, and
$h_{\mathrm{top}}(\widetilde T_\aaa)=0$; a factor of a zero-entropy system has zero entropy.
For the last statement, $\mathrm{Ad}\,\mu_\aaa$ is inner, and inner automorphisms have
zero topological entropy \cite{Brown}.
\end{proof}

\begin{remark}[Why $\log\Ns\aaa$ does not occur]\label{rem:nolog}
The map with entropy $\log\Ns\pp$ is the digit shift $x\mapsto(x-a_0)/\pp$, which is
$\Ns\pp$-to-one; its definition requires \emph{subtracting} the last digit, i.e.\ the
additive structure. The multiplicative system has one isometry per prime, with
range projections $\mathbf 1_{\pp Y_S}$ that are nested rather than orthogonal
($\mathbf 1_{\pp Y_S}\mathbf 1_{\qq Y_S}=\mathbf 1_{\pp\qq Y_S}\ne0$), hence no Cuntz relation and no
$\log n$ (Paper II). The
unavailability of $\log\Ns\aaa$ is the same phenomenon as the failure of the Cuntz relation
there.
\end{remark}

\begin{theorem}[The modular flow has zero entropy]\label{thm:cnt}
Let $\iota:\R\to\widehat{I_S}=\prod_{\pp\in S}\T$, $\iota(t)_\pp=\Ns\pp^{it}$, and let
$\gamma$ be the gauge action. Then $\sigma_t=\gamma_{\iota(t)}$, so $\{\sigma_t\}$ has
compact closure in $\Aut(\cA_{K,S})$; and an automorphism with relatively compact orbit
closure has zero CNT entropy. Hence $h_\varphi(\sigma_t)=0$ for every $\mathrm{KMS}_\beta$
state and every $t$, extremal or symmetric alike.
\end{theorem}

\begin{proof}
$\iota$ is a continuous homomorphism and both $\sigma_t$ and $\gamma_{\iota(t)}$ fix
$C(Y_S)$ pointwise and send $\mu_\pp\mapsto\Ns\pp^{it}\mu_\pp$; the gauge action of the
compact group $\widehat{I_S}$ is continuous for the pointwise norm topology, so its image is
compact. For the entropy, fix a finite-dimensional channel $\gamma$ and $\epsilon>0$; by
relative compactness every $\sigma_{kt}\circ\gamma$ is within $\epsilon$ of one of finitely
many $\theta_j\circ\gamma$, $j\le m(\epsilon)$. Since $H_\varphi$ is symmetric, unchanged by
repetition, subadditive, bounded by $\log\dim$ on a single channel, and continuous with a
modulus independent of the number of arguments,
\[
H_\varphi\bigl(\gamma,\sigma_t\gamma,\dots,\sigma_{(n-1)t}\gamma\bigr)\le m(\epsilon)\log\dim+n\,\delta(\epsilon),
\]
so $\limsup_n\frac1nH_\varphi\le\delta(\epsilon)\to0$. Every KMS state is $\sigma$-invariant,
so the hypothesis holds. The properties of $H_\varphi$ used here are those established in
\cite{CNT}; see also \cite{OP} for the vanishing of $\mathrm{CNT}$ entropy on automorphisms
lying in a compact group.
\end{proof}

\begin{remark}\label{rem:rotation}
This is the noncommutative form of the fact that rotations have zero entropy. Here the
angle is $(\log\Ns\pp)_{\pp\in S}$; over $\Q$ these are $\Q$-linearly independent and
Kronecker's theorem makes $\iota(\R)$ dense in $\widehat{I_S}$, and in general $\iota(\R)$
is dense in the annihilator of $\{\aaa\in I_S:\Ns\aaa=1\}$. The modular flow is an
irrational winding on an infinite torus, dense hence ergodic on gauge orbits, compact hence
entropy-free. Neither property detects
$\Xi_\beta$, which by Corollary~\ref{cor:where} is not an invariant of an automorphism.
\end{remark}

\section{Relative entropy and the information staircase}\label{sec:relent}

Fix $\beta$ with $\Xi_\beta$ finite and write $G=\widehat{\Xi_\beta}$, $A_\beta$ for the set
of thresholds already crossed in the cascade systems of \cite[Prop.~4.16]{Main}, so that
$\Xi_\beta\cong(\Z/2\Z)^{A_\beta}$ and $|G|=2^{|A_\beta|}$.

\begin{theorem}[The staircase]\label{thm:staircase}
Let $\varphi_\varepsilon$, $\varepsilon\in G$, be the extremal $\mathrm{KMS}_\beta$ states
and $\phisym$ the symmetric one. Then
\begin{equation}\label{eq:staircase}
\Ent\bigl(\omega_{\varphi_\varepsilon}\,\big\|\,\omega_{\phisym}\bigr)=\log|G|=|A_\beta|\log2 ,
\end{equation}
computed in $\Msym=\pi_{\phisym}(\cA_{K,S})''=\bigoplus_{\varepsilon}M_\varepsilon$, and the
whole value is carried by the centre: each type $\mathrm{III}$ summand contributes zero. For
$\varepsilon\ne\varepsilon'$, $\Ent(\omega_{\varphi_\varepsilon}\|\omega_{\varphi_{\varepsilon'}})=+\infty$.
When $A_\beta$ is infinite the two states are mutually singular and \eqref{eq:staircase} is
$+\infty$.
\end{theorem}

\begin{proof}
By \cite[Thm.~3.6]{Main} the extreme points of the $\mathrm{KMS}_\beta$ simplex are indexed
by $G$ and their measures are ergodic for the orbit relation and pairwise distinct, hence
mutually singular; so the GNS representations are pairwise disjoint, $\phisym=|G|^{-1}\sum_\varepsilon\varphi_\varepsilon$
gives $\pi_{\phisym}=\bigoplus_\varepsilon\pi_{\varphi_\varepsilon}$ and the central
decomposition $\Msym=\bigoplus_\varepsilon M_\varepsilon$, with central projections
$z_\varepsilon$ of $\omega_{\phisym}$-value $|G|^{-1}$. For a direct sum with $\omega=\bigoplus\lambda_i\omega_i$ and
$\psi=\bigoplus\mu_i\psi_i$ one has the additivity formula
\[
\Ent(\omega\|\psi)=\sum_i\lambda_i\log\frac{\lambda_i}{\mu_i}+\sum_i\lambda_i\,\Ent(\omega_i\|\psi_i),
\]
with the conventions $0\log0=0$ and $\lambda\log(\lambda/0)=+\infty$ for $\lambda>0$; this is
the relative modular operator of a direct sum decomposing as the direct sum of the relative
modular operators scaled by $\lambda_i/\mu_i$ \cite{Araki}, \cite{OP}.
Index the summands by $\eta\in G$. The state $\omega_{\varphi_\varepsilon}$ has
$\lambda_\eta=1$ for $\eta=\varepsilon$ and $\lambda_\eta=0$ otherwise, while
$\omega_{\phisym}$ has $\mu_\eta=|G|^{-1}$ for every $\eta$. The single surviving term of
the first sum is $\log\bigl(1/|G|^{-1}\bigr)=\log|G|$. In the second sum the term
$\eta=\varepsilon$ vanishes because both states induce the \emph{same} normal state of
$M_\varepsilon$, and every other term has $\lambda_\eta=0$; so the second sum is identically
zero, which is the assertion that the relative entropy is purely central. Disjointness gives
$+\infty$ for $\varepsilon\ne\varepsilon'$. When $A_\beta$ is infinite, $\Xi_\beta^\perp$
has infinite index in $G_S$, the extremal measure is Haar measure on a coset of
$\Xi_\beta^\perp$ in each fibre while the symmetric one is Haar measure on $G_S$, so the two
are mutually singular and the relative entropy is $+\infty$.
\end{proof}

\begin{corollary}[One bit per threshold]\label{cor:bit}
Each threshold contributes exactly $\log2$ to \eqref{eq:staircase}, and the function
$\beta\mapsto|A_\beta|\log2$ is a pure jump staircase with jumps $\log2$ at the thresholds.
It coincides with $\log[\Msym:M_\varepsilon]$, the logarithm of the index of the
symmetry-breaking inclusion $M_\varepsilon\cong\Msym^G\subset\Msym\cong M_\varepsilon\otimes\ell^\infty(G)$,
which is $|G|$ (Paper II of the series): the index staircase and the information
staircase are the same object, both counting $\log$ of the number of phases.
\end{corollary}

\begin{remark}[The order parameter]\label{rem:psi}
For the single-threshold system of \cite[Prop.~4.12]{Main} the two phases are separated by
the pointwise invariant $\Psi$, whose almost-everywhere definedness is a Borel--Cantelli
condition. Under $\phisym$ the observable $\Psi$ takes the values $\pm1$ with probability
$\tfrac12$ each, so $H_{\phisym}(\Psi)=\log2$, while $H_{\varphi_\pm}(\Psi)=0$: crossing the
threshold freezes one previously maximally mixed bit. This is \eqref{eq:staircase} in its
elementary form.
\end{remark}

\begin{proposition}[No preferred phase]\label{prop:landauer}
At fixed $\beta$ the $\mathrm{KMS}_\beta$ states form, by \cite[Thm.~3.6]{Main}, the
simplex $\Prob(G_S/\Xi_\beta^\perp)$, on whose extreme points $G_S$ acts transitively. Every
extremal state is therefore carried to every other by a symmetry of the dynamical system,
and $\phisym$ is the barycentre of the $G_S$-orbit of any of them. No quantity attached to
the pair $(\cA_{K,S},\sigma)$ distinguishes one extremal $\mathrm{KMS}_\beta$ state from
another, and none singles out $\phisym$ except its symmetry.
\end{proposition}

\begin{proof}
The action of $G_S$ on the simplex factors through $G_S/\Xi_\beta^\perp$ and permutes the
extreme points simply transitively \cite[Thm.~3.6(3)]{Main}; the symmetries commute with
$\sigma_t$.
\end{proof}

\begin{remark}[What the bit is]\label{rem:landauer}
Landauer's bound concerns \emph{erasure}: an external agent irreversibly resetting a bit,
which requires work $k_BT\log2$. Here nothing is erased and no agent acts. Crossing a
threshold makes an observable well defined almost everywhere, and the system then has two
equilibrium phases in which it takes a definite value. The bit is not destroyed but
\emph{transferred}, from the microscopic description to the label of the phase; the $\log2$
of \eqref{eq:staircase} is the information needed to say which phase one is in, not
information lost. By Proposition~\ref{prop:landauer} no phase is energetically preferred,
so the analogy with Landauer does not survive. (In the Gibbs regime $\zeta_{K,S}(\beta)<\infty$
one can say more: the extremal states are the normalized Gibbs states of \cite[Thm.~3.6]{Main},
all with the same partition function, hence literally the same free energy; in the type
$\mathrm{III}$ regime no free energy is defined, and the symmetry argument is all there is.)
\end{remark}

\section{Assessment}

\[
\begin{array}{lll}
\text{Candidate invariant} & \text{Sees }\Xi_\beta? & \text{Reason}\\\hline
K_*(\cA_{K,S}) & \text{no} & \text{Cor.~\ref{cor:annih}: coinvariants kill the eigenvectors}\\
K_*(\cB_{K,S}) & \text{no} & \text{Thm.~\ref{thm:core}(1)}\\
\text{dual trace on }K_0(\cB_{K,S}) & \text{no} & \text{Thm.~\ref{thm:core}(3): scaling vs.\ homotopy}\\
\hat\sigma_*\ \text{on}\ K_*(\cB_{K,S}) & \text{no} & \text{Thm.~\ref{thm:core}(2): }\R\ \text{connected}\\
\Delta_{\hat\varphi}\ \text{on}\ K_1 & \text{not defined} & \text{Prop.~\ref{prop:dlhs}: domain is }\cU_0\\
H_*(\cG_S;\Z)\ \text{and }c\text{-twists} & \text{no} & \text{Prop.~\ref{prop:homology}: no }\beta\\
h_\varphi(\alpha_\aaa),\ h_\varphi(\cL_\aaa) & \text{not defined} & \text{Thm.~\ref{thm:noinv}: no invariant measure}\\
\htop(\alpha_\aaa) & 0 & \text{Thm.~\ref{thm:ht}: a contraction}\\
h_\varphi(\sigma_t) & 0 & \text{Thm.~\ref{thm:cnt}: compact gauge subflow}\\
\bigl(K_0(C(Y_S)),\alpha_*\bigr) & \textbf{yes} & \text{Thm.~\ref{thm:eigen}}\\
H^1_{\mathrm{meas}}(\cR_S,\pi_\beta;\T) & \textbf{yes} & \text{Thm.~\ref{thm:meas}}\\
\Ent(\omega_{\varphi_\varepsilon}\|\omega_{\phisym}) & \textbf{yes} & \text{Thm.~\ref{thm:staircase}}
\end{array}
\]

The three affirmative entries have a common form: each is built directly from a
Radon--Nikodym derivative. The eigenvalue equation of Theorem~\ref{thm:eigen} is the
statement that $d(\nu\circ\alpha_\pp)/d\nu=\Ns\pp^{-\beta}$; the cohomology of
Theorem~\ref{thm:meas} is computed with respect to the measure class of $\pi_\beta$, that is,
modulo the Radon--Nikodym cocycle; and the relative entropy of
Theorem~\ref{thm:staircase} is by definition
$\int\log\bigl(d\omega/d\psi\bigr)\,d\omega$. Each therefore compares two objects in the
same measure class, rather than iterating a single map or extracting a homotopy invariant of
a single algebra. The negative entries are exactly the constructions that discard the
derivative: a $C^*$-algebra retains no measure, a topological groupoid retains no measure
class, and a homotopy functor cannot see a scaling. That is the content of Corollary~\ref{cor:where}, and it is the reason the
negative entries are uniform.

Three questions remain. We have not computed $H^1_{\mathrm{meas}}(\cR_S,\pi_\beta;\T)$
itself, only its kernel; for the tail relation of a product measure this should be
accessible, and would presumably exhibit $\widehat{G_S}/\Xi_\beta$ inside it. Second,
Lemma~\ref{lem:dict} is stated for $\T$-valued cocycles; the extended group of
\cite{SeriesI}, defined by cocycles valued in $G_S\times\R$, is the corresponding kernel for
the factor type, and there the inclusion $E(\Phi)\subseteq\Lambda^\perp$ of
Lemma~\ref{lem:dict} yields only an upper bound on the type; whether the kernel determines
the type --- i.e.\ whether the cocycle is regular, or equivalently whether the flow of weights
has pure point spectrum --- is open (Remark~\ref{rem:noconverse}). Third, the twisted cyclic $1$-cocycle
of Remark~\ref{rem:parity} is the natural home for a pairing with $K_1(\cB_{K,S})$; whether
the resulting invariant recovers Theorem~\ref{thm:eigen} we do not know.

\begin{thebibliography}{9}
\bibitem{Araki} H. Araki, \emph{Relative entropy of states of von Neumann algebras}, Publ. RIMS Kyoto Univ. 11 (1976), 809--833.
\bibitem{Brown} N. P. Brown, \emph{Topological entropy in exact $C^*$-algebras}, Math. Ann. 314 (1999), 347--367.
\bibitem{CNT} A. Connes, H. Narnhofer, W. Thirring, \emph{Dynamical entropy of $C^*$-algebras and von Neumann algebras}, Comm. Math. Phys. 112 (1987), 691--719.
\bibitem{CT} A. Connes, \emph{An analogue of the Thom isomorphism for crossed products of a $C^*$-algebra by an action of $\R$}, Adv. Math. 39 (1981), 31--55.
\bibitem{CS} A. Connes, E. St\o rmer, \emph{Entropy for automorphisms of $\mathrm{II}_1$ von Neumann algebras}, Acta Math. 134 (1975), 289--306.
\bibitem{dlHS} P. de la Harpe, G. Skandalis, \emph{D\'eterminant associ\'e \`a une trace sur une alg\`ebre de Banach}, Ann. Inst. Fourier 34 (1984), 241--260.
\bibitem{FM} J. Feldman, C. C. Moore, \emph{Ergodic equivalence relations, cohomology, and von Neumann algebras II}, Trans. Amer. Math. Soc. 234 (1977), 325--359.
\bibitem{OP} M. Ohya, D. Petz, \emph{Quantum Entropy and Its Use}, Springer, 1993.
\bibitem{ST} J.-L. Sauvageot, J.-P. Thouvenot, \emph{Une nouvelle d\'efinition de l'entropie dynamique des syst\`emes non commutatifs}, Comm. Math. Phys. 145 (1992), 411--423.
\bibitem{Schmidt} K. Schmidt, \emph{Cocycles on Ergodic Transformation Groups}, Macmillan Lectures in Math. 1, Delhi, 1977.
\bibitem{Voiculescu} D. Voiculescu, \emph{Dynamical approximation entropies and topological entropy in operator algebras}, Comm. Math. Phys. 170 (1995), 249--281.
\bibitem{SeriesI} R. Chen, \emph{Localized Bost--Connes systems I: factor type and the spectrum of transitions}, preprint, 2026, DOI \href{https://doi.org/10.5281/zenodo.22160827}{10.5281/zenodo.22160827}.
\bibitem{Main} R. Chen, \emph{KMS state uniqueness and phase transitions in localized Bost--Connes systems}, preprint, 2026, \newline\url{https://doi.org/10.5281/zenodo.22152101}.
\end{thebibliography}

\end{document}
